\documentclass[10pt,twocolumn]{article}

% ── Packages ──────────────────────────────────────────────────────────────────
\usepackage[margin=1in,columnsep=0.25in]{geometry}
\usepackage{amsmath,amssymb,amsthm}
\usepackage{graphicx}
\usepackage{booktabs}
\usepackage[hidelinks]{hyperref}
\usepackage{microtype}
\usepackage{xcolor}
\usepackage{algorithm}
\usepackage{algpseudocode}
\usepackage{cite}
\usepackage{url}
\usepackage{balance}
\usepackage{parskip}
\usepackage{enumitem}

% tighter list spacing
\setlist{itemsep=2pt,topsep=3pt}

% theorem environments
\newtheorem{theorem}{Theorem}
\newtheorem{lemma}[theorem]{Lemma}
\newtheorem{proposition}[theorem]{Proposition}

% ── Title ─────────────────────────────────────────────────────────────────────
\title{\textbf{GravTree: A Local-Compaction Tree\\
       for Predictable Write Latency}}

\author{
  Subodh Wandile \\
  Independent Researcher, India \\
  \texttt{sawandile\_b17@it.vjti.ac.in}
}

\date{}

% ──────────────────────────────────────────────────────────────────────────────
\begin{document}
\maketitle

% ── Abstract ──────────────────────────────────────────────────────────────────
\begin{abstract}
Log-Structured Merge-trees (LSM-trees), the dominant write-optimised storage
engine, exhibit tail write latency that is difficult to bound because
compaction work is scheduled globally across levels.  We present
\textbf{GravTree}, a tree-structured key-value index in which every node is an
immutable Sorted String Table (SSTable) file and compaction is strictly
\emph{local}: a single compaction event involves exactly one parent node and
its direct children, and nothing else.  This invariant bounds the I/O of any
single compaction event to $O(K \cdot B)$ bytes---a constant independent of
dataset size---so that the maximum work done between any two writes is fixed
by configuration rather than by data volume.

Writes land in an in-memory skip-list at the root; the critical write path
performs no disk I/O.  A dedicated background thread runs LocalFlush to push
data from the root buffer down the tree, splitting leaves and internal nodes
as they fill.  We implement a complete C++17 prototype featuring per-node Bloom
filters, a crash-safe atomic manifest with garbage collection of superseded
nodes, and a two-level lock design that allows concurrent reads during
flushes.

We evaluate the prototype against RocksDB~11.8 at $N = 10$M operations per
workload under equal durability settings (no WAL on either engine).  GravTree
achieves 1.3--2.6$\times$ lower p99 latency on all four workloads and a
predictable tail: p99.99 write latency stays below 20\,$\mu$s, versus
31--2117\,$\mu$s for RocksDB.  Point reads are 2.6$\times$ faster at p99 and
2.4$\times$ higher in throughput.  The price is write throughput: the current
single-threaded flush design rewrites whole nodes on every LocalFlush and
sustains 21--80$\times$ fewer writes per second than RocksDB.  We quantify
this trade-off, and discuss how GravTree's immutable, sequentially-written
node files are structurally compatible with Zoned Namespace (ZNS) SSDs.
\end{abstract}

% ── 1. Introduction ───────────────────────────────────────────────────────────
\section{Introduction}
\label{sec:intro}

Write-optimised storage engines power the data layer of cloud databases,
time-series stores, and messaging systems.  The dominant design,
the Log-Structured Merge-tree (LSM-tree)~\cite{oneil1996lsm}, achieves high
write throughput by buffering writes in memory and periodically merging
sorted runs on disk.  Systems including RocksDB~\cite{rocksdb},
LevelDB~\cite{leveldb}, Cassandra, and TiKV all build on this foundation.

The central cost of the LSM design is \emph{compaction}: a background process
that reads files from level $L_i$, merges them with overlapping files from
level $L_{i+1}$, and writes out a new set of $L_{i+1}$ files.  Compaction
can cascade---a newly enlarged $L_{i+1}$ triggers further compaction into
$L_{i+2}$, and so on---so the I/O cost of a single compaction event is
proportional to the sizes of all affected levels.  When the write rate
outpaces compaction throughput, RocksDB stalls or throttles incoming writes
until compaction catches up.  These stall events are well-documented in
production deployments~\cite{lsm-survey,dostoevsky}: p99.9 write latency
routinely spikes 100$\times$--1000$\times$ above median.

Efforts to tame compaction have focused on reducing \emph{write
amplification}~\cite{wisckey,pebblesdb,dostoevsky,monkey} or restructuring
the level hierarchy~\cite{pebblesdb}, but none bound the \emph{duration} of
a single compaction event to a constant independent of dataset size.

\textbf{GravTree} is built around a single invariant that does exactly this:

\begin{quote}
\textit{A LocalFlush involves exactly one parent node and its direct
children---no other nodes are read or written.}
\end{quote}

Every node in GravTree is an immutable SSTable file.  Incoming writes
accumulate in an in-memory skip list at the root.  When the buffer exceeds a
configurable soft threshold, a background thread runs LocalFlush: it
partitions the buffer by child key ranges, merge-sorts each partition with the
corresponding child's on-disk data, writes new immutable child files, and
invalidates the old ones.  The total I/O per flush is $O(K \cdot B)$ bytes
where $K$ is the fanout ($\leq 8$) and $B$ is the buffer threshold---a bounded
constant.  Writers block only if the buffer reaches a hard limit of twice
the threshold before the flush completes.

GravTree is closest in spirit to the B$^\epsilon$-tree~\cite{bender2007cache}
and its high-performance descendant SplinterDB~\cite{splinterdb}, which also
propagate writes downward through per-node buffers.  The key difference is our
storage model: every GravTree node is a complete, immutable SSTable (not a
mutable fixed-size block).  This enables per-node Bloom filters with no
staleness concern, deterministic crash recovery via atomic file rename, and a
natural mapping onto ZNS SSDs where each node occupies one sequentially-written
zone.

\textbf{Contributions.}
\begin{enumerate}
  \item The GravTree data model and LocalFlush algorithm, with a formal bound
        on per-flush I/O cost and a recursive split rule that keeps every node
        within a configured size as the tree grows (\S\ref{sec:design}).
  \item A complete C++17 implementation with skip-list write buffer, per-node
        Bloom filters, sparse index, CRC-64 integrity, atomic manifest, and
        manifest-driven garbage collection, validated at 10M keys for
        completeness and crash-restart consistency (\S\ref{sec:impl}).
  \item An empirical evaluation at $N = 10$M operations against RocksDB~11.8
        showing 1.3--2.6$\times$ lower p99 latency on all workloads and a
        p99.99 write tail below 20\,$\mu$s, together with an honest
        measurement of the write-throughput and write-amplification cost of
        the design (\S\ref{sec:eval}).
  \item A structural analysis of why GravTree's immutable node model is a
        natural fit for ZNS SSDs (\S\ref{sec:zns}).
\end{enumerate}

% ── 2. Background ─────────────────────────────────────────────────────────────
\section{Background and Motivation}
\label{sec:background}

\subsection{LSM-Trees}

An LSM-tree organises data as a sequence of levels $L_0, L_1, \ldots, L_k$.
Writes accumulate in an in-memory MemTable; when full, the MemTable is flushed
as an immutable SSTable to $L_0$.  When $L_0$ accumulates enough files, a
compaction job merges them with overlapping $L_1$ files to produce new $L_1$
files.  This repeats recursively: a sufficiently large $L_i$ triggers
compaction into $L_{i+1}$.

The total I/O cost of a compaction at level $L_i$ is proportional to the sizes
of all files read and written at levels $L_i$ and $L_{i+1}$, which can be
several GB.  When the write rate exceeds the compaction throughput on the
available I/O bandwidth, RocksDB imposes a write delay (soft limit) or
halts writes entirely (hard limit) until compaction catches up.  The
resulting latency spikes are unbounded in duration---they last as long as
the I/O subsystem needs to drain the backlog.

\subsection{\texorpdfstring{B$^\epsilon$-Trees}{B-epsilon Trees}}

The B$^\epsilon$-tree~\cite{bender2007cache} routes writes as
``messages'' through internal node buffers, batching them before pushing
to children.  A flush operation moves a buffer's worth of messages one
level down, amortising I/O across many logical writes.  The flush scope
is bounded to a single parent-child edge, which is the same locality
property exploited by GravTree.

SplinterDB~\cite{splinterdb} is the state-of-the-art B$^\epsilon$-tree
implementation, achieving write throughput competitive with RocksDB while
offering better tail latency.  GravTree differs from SplinterDB in two ways:
(1) nodes are immutable SSTable files rather than mutable B-tree pages, so
crash recovery requires only an atomic manifest update and no write-ahead log
for tree structure; and (2) per-node Bloom filters are always consistent with
the on-disk content, since a node's entries never change after it is written.

\subsection{The Stall Characterisation Problem}

A key question is: \textit{what bounds the worst-case duration of a single
compaction event?}  For an LSM-tree at level $L_i$ with $S_i$ bytes of data,
the answer is $\Theta(S_i + S_{i+1})$, which grows with the dataset.  For
GravTree, the answer is $O(K \cdot B)$ regardless of dataset size.  This
is the fundamental architectural difference that GravTree exploits.

\subsection{Zoned Namespace SSDs}

ZNS SSDs~\cite{zns} replace the conventional block interface with a set of
\emph{zones} that must be written sequentially and explicitly reset before
reuse.  This eliminates the device-level Flash Translation Layer (FTL) and
its associated garbage collection overhead.  Existing storage engines adapted
for ZNS (e.g., ZenFS~\cite{zenfs}) still inherit the multi-level compaction
behaviour of their host engine.  GravTree's append-only node model aligns
structurally with ZNS requirements without requiring an adaptation layer.

% ── 3. Design ─────────────────────────────────────────────────────────────────
\section{GravTree Design}
\label{sec:design}

\subsection{Data Model}

GravTree is a $K$-ary tree of \emph{nodes}.  Each node $v$ stores:
\begin{itemize}
  \item A non-negative \emph{level} $\ell(v)$; level 0 denotes a leaf.
  \item An immutable \emph{SSTable file} containing a sorted sequence of
        entries $\langle \mathit{key},\, \mathit{value},\, \mathit{seq},\,
        \mathit{deleted} \rangle$.
  \item References to at most $K$ child nodes, each covering a disjoint
        key sub-range.
  \item A Bloom filter and sparse index embedded in the SSTable file.
  \item A \emph{key range} $[\mathit{min\_key},\, \mathit{max\_key}]$
        covering all keys in the subtree rooted at $v$.
\end{itemize}

The \emph{root node} additionally maintains an in-memory skip-list write
buffer that absorbs all incoming writes without disk I/O.  A global
\emph{manifest} file records the set of valid nodes and the current root ID.

\subsection{Write Path}

\begin{algorithm}[t]
\caption{Insert($key$, $value$)}
\label{alg:insert}
\begin{algorithmic}[1]
\If{$\mathrm{buf.bytes}() \geq \mathrm{hardLimit}$}
  \State \Call{WaitForFlush}{} \Comment{back-pressure; rarely taken}
\EndIf
\State $e \gets \langle key,\ value,\ \mathrm{NextSeq}(),\ \mathrm{false} \rangle$
\State $\mathrm{buf.Insert}(e)$ \Comment{$O(\log n)$, no disk I/O}
\If{$\mathrm{buf.bytes}() \geq \mathrm{softThreshold}$}
  \State \Call{SignalFlushThread}{}
\EndIf
\end{algorithmic}
\end{algorithm}

Algorithm~\ref{alg:insert} shows the write path.  The only critical-path work
is a skip-list insert ($O(\log n)$) into DRAM.  A dedicated background thread
is signalled when the buffer exceeds a soft threshold (default 4\,MB); the
thread drains the buffer without blocking writers.  Writers block only at the
hard limit (default $2\times$ soft threshold), which is never reached under
normal sustained write rates.

The skip list orders entries by $(\mathit{key}\ \mathrm{asc},\,
\mathit{seq}\ \mathrm{desc})$, so the most recent version of a key is always
the first encountered during lookup.

\subsection{LocalFlush}

\begin{algorithm}[t]
\caption{LocalFlush(root $r$)}
\label{alg:flush}
\begin{algorithmic}[1]
\State $E \gets \mathrm{buf.ToSortedVector}()$;\ $\mathrm{buf.Clear}()$
\State Partition $E$ into $E_1, \ldots, E_K$ by child key ranges
\For{$i \gets 1$ \textbf{to} $K$}
  \State $E'_i \gets \mathrm{Merge}(E_i,\ \mathrm{Load}(c_i))$ \Comment{two-way merge-sort}
  \State $c'_i \gets \mathrm{WriteSSTable}(E'_i)$ \Comment{new immutable file}
  \State $\mathrm{Invalidate}(c_i)$ \Comment{old file now garbage}
\EndFor
\State Update child pointers in $r$
\State $\mathrm{AtomicManifestSave}()$
\end{algorithmic}
\end{algorithm}

Algorithm~\ref{alg:flush} is the core operation.  Its I/O cost is bounded:

\begin{proposition}
\label{prop:io}
A single LocalFlush reads and writes at most
$B + K \cdot C_{\max}$ bytes, where $B$ is the root buffer size and
$C_{\max}$ is the size of the largest child SSTable.  Both $B$ and
$C_{\max}$ are bounded by configuration constants independent of the
total dataset size $N$.
\end{proposition}

\textbf{Merge semantics.}  Entries from the write buffer and from a child
SSTable are merged by $(\mathit{key}\ \mathrm{asc},\, \mathit{seq}\
\mathrm{desc})$.  For any key, only the entry with the highest
$\mathit{seq}$ is retained in the output.  A tombstone entry
($\mathit{deleted} = \mathrm{true}$) suppresses all earlier versions of
its key.

\subsection{Leaf Split}

When the root is a childless leaf and its buffer reaches the soft threshold, a
leaf split creates two new leaf SSTables and a new internal root:

\begin{enumerate}[leftmargin=1.5em]
  \item Sort the write buffer; choose the median key as the fence key.
  \item Write the left and right halves as two new leaf SSTable files.
  \item Create a new internal root referencing both leaves.
  \item Atomically update the manifest.
\end{enumerate}

The old leaf file is invalidated and can be reclaimed.

\subsection{Read Path}

\begin{algorithm}[t]
\caption{Get($key$)}
\label{alg:get}
\begin{algorithmic}[1]
\If{$e \gets \mathrm{buf.FindKey}(key)$} \Comment{$O(\log n)$, DRAM only}
  \State \textbf{return} $e.\mathrm{deleted}\ ?\ \bot\ :\ e.\mathrm{value}$
\EndIf
\State $\mathit{path} \gets \mathrm{FindPath}(key)$ \Comment{root-to-leaf traversal}
\For{each node $v$ \textbf{in} $\mathit{path}$}
  \If{$v.\mathrm{MightContain}(key)$} \Comment{Bloom filter}
    \If{$e \gets v.\mathrm{BinarySearch}(key)$}
      \State \textbf{return} $e.\mathrm{deleted}\ ?\ \bot\ :\ e.\mathrm{value}$
    \EndIf
  \EndIf
\EndFor
\State \textbf{return} $\bot$
\end{algorithmic}
\end{algorithm}

The read path (Algorithm~\ref{alg:get}) first checks the in-memory skip list
in $O(\log n)$.  If not found, it traverses the root-to-leaf path, using each
node's Bloom filter to skip nodes unlikely to contain the key, then performs a
binary search within the node's sparse-indexed SSTable.

\subsection{Complexity}

Let $n$ be the number of entries in the write buffer, $B$ the buffer threshold
in bytes, $K$ the fanout, and $H = O(\log_K (N/B))$ the tree height for
a dataset of $N$ entries.

\begin{itemize}
  \item \textbf{Write (amortised):} $O(\log n)$ --- skip-list insert, DRAM
        only.
  \item \textbf{LocalFlush I/O:} $O(K \cdot B)$ bytes --- bounded constant
        (Proposition~\ref{prop:io}).
  \item \textbf{Point read:} $O(H \cdot \log n_v)$ disk reads in the worst
        case, where $n_v$ is the number of entries per node; Bloom filters
        reduce the effective constant dramatically in practice.
  \item \textbf{Write amplification:} Each entry is rewritten once per tree
        level it descends through, giving write amplification $\leq H =
        O(\log_K (N/B))$---the same asymptotic bound as a B$^\epsilon$-tree.
\end{itemize}

% ── 4. Implementation ─────────────────────────────────────────────────────────
\section{Implementation}
\label{sec:impl}

The GravTree prototype is approximately 1,500 lines of C++17, open-sourced
under the AGPL license.  We describe the key implementation decisions.

\subsection{Skip-List Write Buffer}

The write buffer is a probabilistic skip list with maximum height 12 and
branching probability $p = 1/4$, matching the parameters used in LevelDB's
MemTable~\cite{leveldb}.  Compared to a sorted \texttt{std::vector}, the skip
list eliminates the $O(n)$ element-shifting cost on insert, which was the
dominant bottleneck in our earlier vector-based prototype and caused a 5$\times$
write latency disadvantage at high buffer occupancy.

\subsection{SSTable File Format}

Each node file is written once and never modified.  Its layout is:

\begin{center}
\small
\begin{tabular}{ll}
\toprule
Section & Size \\
\midrule
Magic number + version & 12 B \\
Header (id, level, num\_children, \ldots) & 64 B \\
Child ID array ($\leq 16$ entries) & 128 B \\
Bloom filter (serialised bit-array) & variable \\
Sorted entries (key, value, seq, deleted) & variable \\
Sparse index (every 128\textsuperscript{th} key + file offset) & variable \\
CRC-64 checksum (covers all preceding bytes) & 8 B \\
\bottomrule
\end{tabular}
\end{center}

Files are written atomically via a temporary file renamed into place, ensuring
that a partially-written file is never visible as a valid node.

\subsection{Bloom Filters}

Each on-disk node carries a Bloom filter targeting a 1\% false-positive rate,
built using MurmurHash3-128 with the Kirsch-Mitzenmacher double-hashing
technique~\cite{bloom}.  Filters are serialised into the SSTable file and
loaded lazily on first access.

One subtlety: the root node's in-memory write buffer grows after the Bloom
filter was last serialised, so the filter may not reflect recently buffered
keys.  We handle this by bypassing the Bloom check for any node that has a
non-empty in-memory buffer, falling back to a full binary search.  Once a node
is flushed and its buffer is empty, the Bloom filter is authoritative.

\subsection{Manifest and Crash Safety}

A text-format manifest records: the set of valid node files with their IDs,
levels, and paths; the current root ID; and a monotonically increasing
next-node-ID counter.  All manifest mutations are performed via write-to-temp
+ rename, so the manifest is always either the old or new state---never
partially updated.

On restart, GravTree reads the manifest, ignores any node file whose CRC-64
check fails, and loads the root node into the in-memory node cache.

\textbf{Durability limitation.}  The current prototype does not maintain a
write-ahead log (WAL) for the in-memory write buffer.  A crash before
\texttt{flush()} is called will lose any writes buffered since the last flush.
Adding a WAL or making the skip list itself persistent (e.g., via NVM) is
straightforward future work.

\subsection{Concurrency Model}

GravTree uses a two-level lock hierarchy to separate the hot write path from
the tree-structure path:

\begin{itemize}
  \item \textbf{\texttt{write\_buf\_mu}} (\texttt{std::mutex}): guards the
        skip-list write buffer.  Held only for the duration of a skip-list
        insert---$O(\log n)$ DRAM operations, no disk I/O.  This is the only
        lock on the critical write path.
  \item \textbf{\texttt{cache\_mu}} (\texttt{std::shared\_mutex}): guards the
        node cache and tree structure.  The background flush thread acquires
        an exclusive lock while modifying tree structure; concurrent readers
        (\texttt{get}, \texttt{range\_scan}) hold a shared lock, permitting
        parallel reads.
\end{itemize}

The two locks are never held simultaneously, ruling out deadlock.  The flush
thread releases \texttt{cache\_mu} before writing SSTable files to disk, so
disk I/O does not block readers.

% ── 5. Evaluation ─────────────────────────────────────────────────────────────
\section{Evaluation}
\label{sec:eval}

\subsection{Experimental Setup}

\textbf{Hardware.}  All experiments run on a Linux x86\_64 virtual machine
with 8 vCPUs (Intel Xeon Platinum 8559C), 31\,GB RAM, no swap, and a 124\,GB
NVMe-backed virtual disk (ext4), running Ubuntu 22.04.5 LTS.  Both engines
write to the same filesystem.

\textbf{Software.}  GravTree is built with GCC in Release mode
(\texttt{-O3}).  We compare against RocksDB~11.8.1 built from source with the
same compiler.

\textbf{Configuration.}  Both engines use a 4\,MB write buffer.  RocksDB is
configured with \texttt{kNoCompression}, \texttt{max\_write\_buffer\_number}
$= 3$, one background compaction thread and one background flush thread.
GravTree uses fanout $K = 8$ (\texttt{MAX\_FANOUT} $= 16$ for internal nodes)
and hard-limit factor 2.0.  \emph{Durability is equalised}: the GravTree
prototype has no write-ahead log, so RocksDB is run with
\texttt{disableWAL = true}.  Both engines therefore offer the same guarantee
(buffered writes are lost on crash) and neither pays a per-write log
\texttt{fsync}.

\textbf{Methodology.}  Each workload performs $N = 10{,}000{,}000$ operations
with 100-byte values and 20-byte zero-padded decimal keys.  Latency is
measured per-operation with \texttt{std::chrono::high\_resolution\_clock};
results are reported as sorted percentiles with no warm-up exclusion.
Throughput is wall-clock operations per second over the timed loop
\emph{including} the closing \texttt{flush()}, so it reflects the total cost
of making all $N$ writes persistent.  Before benchmarking, a separate
verification harness inserted 1M keys, read every key back, closed the
engine, re-opened it from the manifest, and read every key again; all reads
returned the correct value (0 missing, 0 mismatched).

\subsection{Workloads}

\begin{itemize}
  \item \textbf{SeqWrite}: Keys $0, 1, \ldots, N{-}1$ in order.
  \item \textbf{RandWrite}: Keys drawn uniformly at random from $[0, 2N)$.
  \item \textbf{PointRead}: Dataset pre-filled with $N/10 = 1$M sequential
        keys ($\approx 120$\,MB, 30$\times$ the write buffer, so most reads go
        to disk); $N$ random point lookups.
  \item \textbf{Mixed}: $N/2$ sequential keys pre-inserted; $N$ operations
        alternating random reads and appending writes.
\end{itemize}

\subsection{Write Latency}

\begin{table}[t]
\centering
\footnotesize
\caption{Write latency ($\mu$s) and throughput at $N = 10$M.
\textbf{Bold} = lower latency.  Ratio $>1$ favours GravTree.}
\label{tab:write}
\begin{tabular}{lrrrrr}
\toprule
 & p50 & p99 & p99.9 & p99.99 & Mops/s \\
\midrule
\multicolumn{6}{l}{\textit{SeqWrite}} \\
GravTree & \textbf{0.26} & \textbf{0.80} & \textbf{1.72} & \textbf{9.2} & 0.067 \\
RocksDB  & 0.62 & 1.28 & 5.45 & 31.4 & \textbf{1.439} \\
Ratio    & 2.4$\times$ & 1.6$\times$ & 3.2$\times$ & 3.4$\times$ & 0.047$\times$ \\
\midrule
\multicolumn{6}{l}{\textit{RandWrite}} \\
GravTree & \textbf{0.72} & \textbf{2.43} & \textbf{7.21} & \textbf{17.3} & 0.003 \\
RocksDB  & 1.08 & 3.20 & 1061.6 & 2116.7 & \textbf{0.232} \\
Ratio    & 1.5$\times$ & 1.3$\times$ & 147$\times$$^\ddagger$ & 122$\times$$^\ddagger$ & 0.013$\times$ \\
\midrule
\multicolumn{6}{l}{\textit{Mixed (50/50)}} \\
GravTree & \textbf{1.54} & \textbf{4.33} & \textbf{9.91} & \textbf{19.2} & 0.103 \\
RocksDB  & 2.43 & 9.81 & 17.25 & 83.0 & \textbf{0.309} \\
Ratio    & 1.6$\times$ & 2.3$\times$ & 1.7$\times$ & 4.3$\times$ & 0.33$\times$ \\
\bottomrule
\multicolumn{6}{p{0.95\linewidth}}{{\footnotesize $^\ddagger$RocksDB's RandWrite
p99.9/p99.99 are dominated by a single write stall (1--4\,ms per operation)
confined to the final $\approx$40K operations; see \S\ref{sec:timeseries}.}}
\end{tabular}
\end{table}

Table~\ref{tab:write} shows write latency and throughput.  At p50 GravTree is
1.5--2.4$\times$ faster: a write is a skip-list insert into DRAM
($\approx 0.26\,\mu$s), whereas RocksDB's memtable insert plus write-group
coordination costs $\approx 0.6$--$1.1\,\mu$s.  With the WAL disabled on both
engines this gap is purely in-memory overhead, not durability cost.

At p99, GravTree's advantage is 1.3--2.3$\times$.  At p99.9 and p99.99 the
picture is qualitatively different: GravTree's tail stays below 20\,$\mu$s on
every write workload, whereas RocksDB's p99.99 ranges from 31\,$\mu$s
(SeqWrite) to 2.1\,ms (RandWrite).  The RandWrite outlier is a RocksDB write
stall, not steady-state behaviour: with a 4\,MB memtable and
\texttt{max\_write\_buffer\_number} $= 3$, the writer eventually outruns the
single flush thread and is throttled.  GravTree's writer is throttled too
(the hard limit at twice the buffer threshold blocks \texttt{insert()} until LocalFlush completes),
but that back-pressure is exercised rarely at this scale and its cost lands
on a small number of operations.

\textbf{Throughput.}  The right-hand column is the cost of the design.
RocksDB sustains 21$\times$ (SeqWrite) to 80$\times$ (RandWrite) more writes
per second.  Two factors dominate.  First, every LocalFlush rewrites whole
nodes: a 4\,MB root flush that touches $c$ children rewrites up to
$c \cdot B$ bytes, and when a child's own buffer fills it recursively rewrites
its children.  For uniformly random keys every root flush touches all
children, so write amplification grows with the number of tree levels and
the fanout.  We measured it directly (bytes passed to \texttt{write(2)} by
the process, from \texttt{/proc/self/io}, divided by the $N \cdot 120$\,B
logical payload; background compactions drained before reading the
counter).  At 1M operations GravTree writes $6.5\times$ (SeqWrite) and
$8.6\times$ (RandWrite) the logical volume versus $1.0\times$ and
$3.3\times$ for RocksDB.  At 10M operations SeqWrite stays at $9.4\times$
(RocksDB $1.0\times$), but RandWrite reaches $\mathbf{180\times}$
(216\,GB written for 1.2\,GB of payload) versus $10.0\times$ for RocksDB:
with uniformly random keys every root flush is spread across all children,
so each flush rewrites essentially the whole tree below it, and the rewrite
volume grows with the dataset size.  This, not per-operation
overhead, is why the 10M RandWrite run took 51 minutes.  Second, the flush thread is single-threaded and performs merge,
Bloom-filter construction, CRC-64 and file write serially while holding the
tree lock.  The single-threaded flush is a prototype limitation; the random-write
amplification is a property of whole-node rewriting under the current
buffer/fanout configuration and is the limitation a production deployment
would hit first (\S\ref{sec:limits}).  Reducing it requires either
larger child buffers relative to the root (so a flush touches fewer
children) or partial-node rewrites, both left to future work.

\subsection{Read Latency}

\begin{table}[t]
\centering
\small
\caption{Point read latency ($\mu$s) and throughput, 1M-key dataset,
10M random lookups.  \textbf{Bold} = better.}
\label{tab:read}
\begin{tabular}{lrrrrr}
\toprule
 & p50 & p99 & p99.9 & p99.99 & Mops/s \\
\midrule
GravTree & \textbf{1.47} & \textbf{2.42} & \textbf{9.29} & \textbf{14.4} & \textbf{0.620} \\
RocksDB  & 3.97 & 6.21 & 13.39 & 50.7 & 0.258 \\
Ratio    & 2.7$\times$ & 2.6$\times$ & 1.4$\times$ & 3.5$\times$ & 2.4$\times$ \\
\bottomrule
\end{tabular}
\end{table}

Point reads are GravTree's strongest result (Table~\ref{tab:read}):
2.6$\times$ lower p99, 3.5$\times$ lower p99.99, and 2.4$\times$ higher
throughput on a dataset 30$\times$ larger than the write buffer, so the
result is not a DRAM-hit artefact.  A GravTree lookup walks a tree of height
2--3, consulting one Bloom filter per node on the path and reading at most
one leaf; there is no level-by-level probing of overlapping SSTables and no
block cache indirection.

\subsection{Latency Time-Series}
\label{sec:timeseries}

\begin{table}[t]
\centering
\footnotesize
\caption{Peak write latency per 20\,000-operation window (500 windows per
workload), all in $\mu$s: range of per-window peaks over the run, count of
windows whose peak exceeds 100\,$\mu$s, and the single worst window.  Both engines'
final-window spike is the explicit \texttt{flush()} issued after the last
operation and is excluded.  $^\ddagger$ops 9.96M--9.98M.}
\label{tab:timeseries}
\setlength{\tabcolsep}{3pt}
\begin{tabular}{llrrr}
\toprule
Workload & Engine & Range & $>$100 & Worst \\
\midrule
SeqWrite  & GravTree & 0.2--1.4  & 0 & 1.4 \\
          & RocksDB  & 0.6--2.7  & 0 & 2.7 \\
RandWrite & GravTree & 0.3--4.2  & 0 & 4.2 \\
          & RocksDB  & 0.8--6.7  & 1 & 1\,060$^\ddagger$ \\
Mixed     & GravTree & 0.2--8.7  & 0 & 8.7 \\
          & RocksDB  & 0.6--15.3 & 0 & 15.3 \\
\bottomrule
\end{tabular}
\end{table}

Table~\ref{tab:timeseries} summarises 500 per-window peaks per workload.
GravTree exhibits no mid-run window above 100\,$\mu$s on any workload: its
peak stays within 0.2--1.4\,$\mu$s for SeqWrite and rises smoothly from
0.3 to 4.2\,$\mu$s over the RandWrite run as leaves fill and the tree
deepens.
RocksDB is likewise flat for SeqWrite and Mixed at this scale; the
\emph{multi-level compaction cascades} that produce tens-of-millisecond
stalls in production did not manifest on this hardware at $N = 10$M with
1.2\,GB of data, and we do not claim to have observed them.  The one RocksDB
excursion is a write stall of $\approx$1\,ms in the final pre-flush RandWrite
window (ops 9.96M--9.98M),
which accounts for its p99.9/p99.99 in Table~\ref{tab:write}.

The final-window spike (excluded above) is the benchmark's closing
\texttt{flush()}.  For GravTree it drains the root buffer through every
level of the tree in one synchronous call---1.2\,s (SeqWrite) to 29.7\,s
(RandWrite)---and is the clearest single measurement of the write
amplification discussed above.  For RocksDB it is a memtable flush of
4--11\,ms.  Neither is an operational stall, but the GravTree figure shows
that a bounded \emph{per-event} cost does not imply a small \emph{total}
cost when many events are chained.

\subsection{Discussion and Limitations}
\label{sec:limits}

\textbf{Write throughput and amplification.}  The dominant limitation of the
current prototype is write throughput: 0.067\,Mops/s sequential and
0.003\,Mops/s random, versus 1.44 and 0.23 for RocksDB, driven by a
measured write amplification of 180$\times$ at 10M random keys (9.4$\times$
sequential) against 10$\times$ and 1.0$\times$ for RocksDB.  The LocalFlush
bound guarantees that no single compaction event exceeds $O(K \cdot B)$ I/O,
but it does not reduce total I/O: each key is rewritten once per tree level
plus once per LocalFlush that touches its leaf, and the flush thread executes
these rewrites serially under the tree lock, so a writer that outruns it is
blocked at the hard limit.  Three engineering changes, none of which alter
the LocalFlush invariant, would close most of this gap and are the immediate
next step: (i) parallel LocalFlush across disjoint children, which are
independent by construction; (ii) delta files for internal nodes, so a flush
into a child appends a run rather than rewriting the node; and (iii) larger
leaves with a smaller root buffer, trading a modest increase in per-event
bound for far fewer rewrites.  Until these are implemented, GravTree is
suited to read-heavy or latency-sensitive workloads with moderate write
rates, not to bulk ingest.

\textbf{What the evaluation does and does not show.}  It shows that
GravTree's write tail is tightly bounded (p99.99 $< 20\,\mu$s on every
workload) and that its p99 is lower than RocksDB's on every workload.  It
does \emph{not} show RocksDB compaction cascades: at 10M operations and
1.2\,GB on a server-class NVMe device, RocksDB's mid-run peaks stayed under
100\,$\mu$s.  Larger datasets, larger values, or slower devices would be
needed to test the ``stall elimination'' hypothesis directly, and we leave
that as an explicit open question rather than a claim.

\textbf{No WAL; durability parity.}  The prototype does not persist the
write buffer to a WAL, and RocksDB was run with \texttt{disableWAL} to match.
Absolute RocksDB write latencies with its default WAL-enabled configuration
are higher (its p50 was $\approx 3\,\mu$s in our original WAL-enabled 100K measurements),
so the comparison here is conservative toward RocksDB.  Adding a WAL to
GravTree is orthogonal to LocalFlush but would add the same per-write cost.

\textbf{Single-threaded writers.}  Concurrent external writers are serialised
by \texttt{write\_buf\_mu}; the numbers above are single-client.  Per-shard
skip lists are straightforward future work.

\textbf{Benchmark platform.}  Measurements were taken on a cloud VM with a
virtualised NVMe disk; absolute values will differ on bare metal, and the
relative results should be re-validated there.

% ── 6. ZNS Compatibility ──────────────────────────────────────────────────────
\section{ZNS SSD Compatibility}
\label{sec:zns}

Zoned Namespace (ZNS) SSDs~\cite{zns} eliminate the device-level FTL by
requiring applications to write each zone sequentially and reset zones before
reuse.  This shifts GC responsibility to the host software.  Storage engines
designed for random-overwrite block devices require substantial rearchitecting
for ZNS compatibility.  ZenFS~\cite{zenfs} adds ZNS support to RocksDB by
mapping SST files to zones, but compaction still rewrites multi-gigabyte zones
non-sequentially within the RocksDB level hierarchy.

GravTree's immutable node model aligns with ZNS requirements by construction:

\begin{itemize}
  \item \textbf{Sequential write.}  Each SSTable file is written front-to-back
        in a single pass (magic, header, Bloom filter, entries, sparse index,
        CRC) and never modified.  This matches the ZNS per-zone sequential
        write requirement exactly.
  \item \textbf{Natural GC.}  When LocalFlush creates new child files, the old
        child files are atomically invalidated.  The corresponding zones can be
        immediately reset and recycled, providing deterministic, bounded GC
        with no separate GC thread.
  \item \textbf{Bounded zone pressure.}  A single LocalFlush activates at most
        $K + 1$ zones simultaneously (one per new child file, plus the parent).
        This simplifies zone allocation compared to LSM compaction, which may
        open $O(|L_i| + |L_{i+1}|)$ files concurrently.
\end{itemize}

A ZNS backend using the libzbd API~\cite{libzbd} is in development.  We plan
to report quantitative results on ZNS hardware in a follow-up evaluation.

% ── 7. Related Work ───────────────────────────────────────────────────────────
\section{Related Work}
\label{sec:related}

\textbf{LSM-tree optimisations.}  A large body of work optimises LSM-tree
compaction policy without bounding its scope.  Dostoevsky~\cite{dostoevsky}
and Monkey~\cite{monkey} tune Bloom filter allocations across levels for
different read/write trade-offs.  WiscKey~\cite{wisckey} separates keys from
values to reduce compaction I/O.  PebblesDB~\cite{pebblesdb} introduces guards
to partition levels and reduce write amplification.  None of these approaches
bound per-compaction I/O to a constant independent of dataset size.

\textbf{\texorpdfstring{B$^\epsilon$-trees.}{B-epsilon trees.}}  Bender et al.~\cite{bender2007cache} introduced
the \texorpdfstring{B$^\epsilon$-tree}{B-epsilon-tree}; TokuDB (now Percona TokuMX) commercialised it.
SplinterDB~\cite{splinterdb} achieves high throughput with small per-flush
scopes.  GravTree uses the same parent-local flush insight but commits to
immutable SSTable files as nodes, enabling consistent per-node Bloom filters
and direct ZNS zone mapping without adaptation.

\textbf{NVMe-optimised LSM-trees.}  SpanDB~\cite{spandb} places the WAL and
top LSM levels on fast NVMe, reducing write latency.  These are orthogonal
to GravTree's compaction-locality property and could be combined.

\textbf{ZNS-aware storage.}  ZenFS~\cite{zenfs} adapts RocksDB for ZNS.
F2FS and dm-zoned provide ZNS-friendly file systems.  None structurally
guarantee that a single compaction touches a bounded number of zones;
GravTree does.

% ── 8. Conclusion ─────────────────────────────────────────────────────────────
\section{Conclusion and Future Work}
\label{sec:conclusion}

We presented GravTree, a tree-structured key-value store whose central
invariant---LocalFlush touches exactly one parent and its direct
children---bounds per-compaction I/O to a constant independent of dataset
size.  A C++17 prototype evaluated at 10M operations against RocksDB under
equal (no-WAL) durability achieves 1.3--2.6$\times$ lower p99 latency on all
four workloads, keeps p99.99 write latency below 20\,$\mu$s, and delivers
2.4$\times$ higher point-read throughput.  The cost is write throughput:
whole-node rewrites give a measured write amplification of 9.4$\times$
(sequential) and 180$\times$ (random) at 10M keys versus 1.0$\times$ and
10$\times$ for RocksDB, and, with a single-threaded flush, a 21--80$\times$
lower sustained write rate.  We did not observe LSM
compaction cascades in RocksDB at this scale, so bounded-cascade behaviour
remains an architectural argument rather than an empirical one.

Planned future work:

\begin{enumerate}[leftmargin=1.5em]
  \item \textbf{Write-amplification reduction.} Partial-node rewrites or
        level-dependent buffer sizes so a random-key flush does not rewrite
        the whole subtree; parallel flush of sibling children with merge,
        Bloom-filter construction and CRC-64 moved off the tree lock.
  \item \textbf{Write-ahead log.} Adding crash durability for the in-memory
        write buffer.
  \item \textbf{ZNS backend.} Quantitative evaluation on ZNS hardware with
        the libzbd backend.
  \item \textbf{Bare-metal and slow-device revalidation.} Repeating the
        evaluation on dedicated hardware and on devices where LSM compaction
        is I/O-bound, where cascade stalls are expected to appear.
  \item \textbf{Formal analysis.} Proof of the write-amplification bound
        and comparison with the B$^\epsilon$-tree optimality result of
        Bender et al.
\end{enumerate}

The GravTree prototype is available as open-source software under the AGPL
license.

% ── References ────────────────────────────────────────────────────────────────
\balance
\bibliographystyle{plain}
\bibliography{refs}

\end{document}
